\documentclass[12pt,a4paper]{article}

\usepackage[utf8]{inputenc}
\usepackage[T1]{fontenc}
\usepackage{geometry}
\usepackage{booktabs}
\usepackage{array}
\usepackage{hyperref}
\usepackage{parskip}
\usepackage{microtype}
\usepackage[backend=biber,style=numeric,sorting=none]{biblatex}

\geometry{margin=2.5cm}

\addbibresource{../report/references.bib}

\hypersetup{
    colorlinks=true,
    linkcolor=black,
    urlcolor=blue,
    citecolor=black
}

\title{\textbf{Bioinformatics Master's Project Plan}}
\date{}

\begin{document}

\maketitle

\section*{General Information}

\begin{tabular}{@{}ll}
    \textbf{Name of student:}      & Oliver E. Todreas \\
    \textbf{Contact email:}        & \href{mailto:ol6437ek-s@student.lu.se}{ol6437ek-s@student.lu.se} \\
    \textbf{Name of supervisor:}   & Parashar Dhapola \\
    \textbf{Department/Host:}      & Nygen Analytics AB \\
    \textbf{Contact email:}        & \href{mailto:parashar@nygen.io}{parashar@nygen.io} \\
    \textbf{Course code:}          & BINP39 \\
    \textbf{Start–end date:}    & March 23, 2026–August 28, 2026 \\
\end{tabular}

\section*{Abstract}

Disease-focused single-cell atlases will be built from SC Base Count, a large standardized scRNA-seq resource that currently lacks systematic clustering and annotation. Disease priorities will be informed by a structured survey of CELLxGENE Census metadata. The work implements a reproducible pipeline (scVI, scanpy, CyteType; AnnData; GitHub) to curate 2–3 disease areas, process datasets individually over an extended per-dataset phase, integrate batches, and benchmark scGPT embeddings against scVI-, PCA/Harmony-, and STATE (Arc Institute) annotation workflows. Evaluation combines concordance metrics (e.g.\ ARI, NMI) with documented edge cases for manuscript-ready comparison.

\section*{Project Plan}

\subsection*{Introduction}

Aggregating many scRNA-seq studies requires consistent preprocessing, batch integration, and interpretable cell typing \cite{placeholder1, placeholder2}. scBaseCount standardizes data at scale but does not yet provide systematic clustering and annotation, limiting cross-cohort disease comparison. Nygen's CyteType delivers evidence-traced, ontology-linked annotations suitable for auditing at scale. How variational integration (scVI)\cite{lopezDeepGenerativeModeling2018}, linear reductions with batch correction (scanpy/Harmony) \cite{placeholder1}, and foundation embeddings (scGPT) \cite{cuiScGPTBuildingFoundation2024} interact with CyteType on this resource is an open, testable question \cite{ahujaMultiagentAIEnables2025}. Candidate disease areas will be evaluated using CELLxGENE Census metadata (tissues, dataset counts, cells, available fields) before joint selection of 2--3 priorities aligned with the company's therapeutic focus.
Aggregating many scRNA-seq studies requires consistent preprocessing, batch integration, and interpretable cell typing \cite{placeholder1, placeholder2}. scBaseCount standardizes data at scale but does not yet provide systematic clustering and annotation, limiting cross-cohort disease comparison. Nygen's CyteType delivers evidence-traced, ontology-linked annotations suitable for auditing at scale. How variational integration (scVI) \cite{lopezDeepGenerativeModeling2018}, linear reductions with batch correction (scanpy/Harmony) \cite{placeholder1}, and foundation embeddings (scGPT) \cite{cuiScGPTBuildingFoundation2024} interact with CyteType on this resource is an open, testable question \cite{placeholder5}. Candidate disease areas will be evaluated using CELLxGENE Census metadata (tissues, dataset counts, cells, available fields) before joint selection of 2–3 priorities aligned with the company's therapeutic focus.

\subsection*{Hypothesis and Specific Aims}

\textbf{H1:} For at least one prioritized disease area, scVI integration will yield CyteType label assignments that differ from per-dataset processing in a pattern measurable by concordance metrics (e.g.\ ARI, NMI); whether integration improves or harms annotation consistency will be assessed per dataset and subset.

\textbf{H2:} CyteType on scGPT-derived structures will not be identical to scVI- or PCA/Harmony-based workflows on the same cells; concordance and discordance versus STATE (Arc Institute) will be quantified and interpreted for biological plausibility.

\textbf{Aims:}
\begin{itemize}
    \item[\textbf{(A)}] Survey CELLxGENE Census metadata; assess 5–6 candidate disease areas; jointly select 2–3 priority areas and subset SC Base Count accordingly.
    \item[\textbf{(B)}] Per dataset: scVI, multi-resolution clustering, subclustering as needed, CyteType, documented quality.
    \item[\textbf{(C)}] Integrate per disease area; compare integrated versus individual CyteType outputs.
    \item[\textbf{(D)}] Benchmark scGPT + CyteType against the above and STATE in a structured four-way comparison.
\end{itemize}

\subsection*{Methods}

\textbf{Data:} SC Base Count disease subsets; AnnData; parameterized pipelines with logging and version control. Phase 1 uses CELLxGENE Census metadata to shortlist disease contexts with sufficient datasets and cells.

\textbf{Procedure:} Phases follow the extended host specification:
\begin{itemize}
    \item \textbf{Phase 1 (weeks 1–3):} Disease atlas curation and selection.
    \item \textbf{Phase 2 (weeks 4–10):} Per-dataset scVI, clustering/subclustering, manual cluster QC, CyteType, edge-case logs.
    \item \textbf{Phase 3 (weeks 11–16):} scVI integration per disease, CyteType on integrated objects, comparison to per-dataset annotations, manuscript-style notes.
    \item \textbf{Phase 4 (weeks 17–20):} scGPT embeddings and CyteType versus scVI-based (Phases 2–3), scanpy PCA/Harmony + CyteType, and STATE (Arc Institute).
\end{itemize}
Tools: Python, scanpy, scvi-tools, scGPT, CyteType SDK.

\textbf{Evaluation:} Hypotheses are tested with (i) quantitative agreement between partitions or label vectors where cells are comparable; (ii) qualitative review of discordant clusters.
\begin{itemize}
    \item \textbf{Support H1:} sustained cross-batch coherence of labels after integration without loss of disease-relevant separation.
    \item \textbf{Refute H1:} integration collapses biologically distinct states or degrades metrics without compensating batch alignment.
    \item \textbf{Support H2:} non-random, interpretable differences across embedding sources at matched resolution.
    \item \textbf{Refute H2:} near-identical outputs across scGPT, scVI, and PCA/Harmony at chosen clustering depth.
\end{itemize}
The pipeline serves these biological comparisons; success is defined by evaluable metrics and documented interpretation, not by tool deployment alone.

\subsection*{Time Plan}

\begin{table}[h!]
\centering
\begin{tabular}{>{\bfseries}p{1.5cm} p{1.8cm} p{10cm}}
\toprule
Weeks & Phase & Activities \\
\midrule
1–3  & Phase 1 & Literature; CELLxGENE Census metadata survey; structured assessment of 5–6 candidate disease areas (tissues, dataset counts, cells, metadata); joint selection of 2–3 priorities. \\[6pt]
4–10 & Phase 2 & Per-dataset scVI training; multi-resolution clustering and subclustering where warranted; manual cluster inspection; CyteType; document annotation quality, edge cases, and CyteType improvement areas. \\[6pt]
11–16 & Phase 3 & scVI batch integration per disease area; CyteType on integrated data; compare with individual-dataset annotations; evaluate where integration helps or hurts; manuscript-ready documentation. \\[6pt]
17–20 & Phase 4 & scGPT embeddings on constructed atlases; CyteType on scGPT-derived clusters; systematic comparison to scVI-, PCA/Harmony-, and STATE-based workflows; structured comparison for manuscript
\bottomrule
\end{tabular}
\end{table}

\textbf{Course deliverables:} Reserve time in the final project weeks for synthesis, final report drafting, and GitHub + README polish; submit the \textbf{final report and repository one week before the seminar presentation} (per course guidelines). Weekly progress reviews with the supervisor continue throughout; remote or in-office work; optional Nygen Lund office days; cloud, compute, and LLM resources as provided by the host.

\printbibliography

\end{document}
